\uuid{r5ft}
\titre{Primitives de fonctions }

\niveau{} 				%L1, L2, L3, MPSI, MP, PCSI, PC, PSI...
\module{ Intégration } 	%Analyse, Algèbre...
\chapitre{Primitive et intégration}   			%Continuité, Groupes, Fonctions de plusieurs variables...
\sousChapitre{}			%Optimisation, Diagonalisation d'une matrice, Calcul de dérivées partielles...

\theme{}				%Fonction de répartition, Division euclidienne de polynômes, ...
\auteur{Quentin Liard}
\datecreate{ 2026-02-11}
\organisation{AMSCC}			%AMSCC, Exo7, ...
\difficulte{1}			%1, 2, 3, 4 ou 5

\contenu{

\texte{ 

}

\begin{enumerate}
\item   \question{Déterminer toutes les primitives de $f(x)=\dfrac{e^{3x}}{1+e^{3x}}$.}
\indication{}
\reponse{Posons $u=1+e^{3x}$, alors $du=3e^{3x}\,dx$. Donc $\displaystyle \int \frac{e^{3x}}{1+e^{3x}}dx=\frac13\int \frac{du}{u}=\frac13\ln(1+e^{3x})+C$.}

\item   \question{Déterminer toutes les primitives de $g(x)=\cos x\,\sin^2 x$.}
\indication{}
\reponse{Avec $u=\sin x$, $du=\cos x\,dx$. Alors $\displaystyle \int \cos x\,\sin^2 x\,dx=\int u^2\,du=\frac{u^3}{3}+C=\frac{\sin^3 x}{3}+C$.}

\item   \question{Déterminer toutes les primitives de $h(x)=\dfrac{3x}{\sqrt{1+x^2}}$.}
\indication{}
\reponse{Posons $u=1+x^2$, $du=2x\,dx$. Alors $\displaystyle \int \frac{3x}{\sqrt{1+x^2}}dx=\frac32\int u^{-1/2}\,du=3\sqrt{u}+C=3\sqrt{1+x^2}+C$.}

\end{enumerate}

}